\scanpage{66}\mspage{7}
% archive sheet number: 68

coïncident. Cela signifie donc que $X$ est déduit de $X^*$ par ``bouchage de trous \unsure{pointés|ponctués} canonique'' (utilisant les structures toriques sur les comp. conn. de $\partial X^*$). \deleted{La question} \inserted{Disons alors} qu'un système de disques $(D_i)_{i\in I}$ dans $X$ \deleted{\unsure{qui se}} autour des $a_i$ est ``admissible''.

\uline{Questions} Comment cette condition d'admissibilité se reflète-t-elle en termes de la géométrie interne de $X$? Par exemple, \deleted{quand} quand $X$ est muni d'une métrique de Bergman canonique (cas où $X\not\simeq\mathbb P^1_{\mathbb C}$), quelle est la relation de cette condition avec celle que les \deleted{$\partial$}$D_i$ soient des ``disques géodésiques'' (\inserted{définis} par la condition que $\text{dist. géod}(x,a_i)\leq r_i$)?

% slanted
\margin{Supposons $(g,\nu)$ hyperbolique i.e. $\ne(0,1),(0,2)$}
Fixons une \deleted{$\unsure{\Sigma_0|X_0}$} surface \unsure{top|tp} \inserted{orientée} à bord de référence \deleted{\unsure{(}} $X_0$, de type $g,\nu$,
